\chapter{Artificial Intelligence Use Declaration}
\label{app:ai_declaration}

This project was developed with assistance from AI-based development tools,
consistent with the University of Hull's Student AI Use Policy (2025).
This appendix provides a full account of the scope of that assistance,
distinguishing clearly between AI-assisted implementation and the candidate's
independent intellectual contributions.

% ─────────────────────────────────────────────────────────────────────────────
\section*{AI-Assisted Components}
\label{app:ai_assisted}
% ─────────────────────────────────────────────────────────────────────────────

AI assistance (Claude, Anthropic) was employed for the following
implementation tasks:

\begin{itemize}[leftmargin=2em]
  \item \textbf{Infrastructure scaffolding.} Supabase project setup, Row-Level
    Security policy boilerplate, Resend email integration for password reset,
    and bcrypt helper functions.
  \item \textbf{Frontend component implementation.} UI components implemented
    from candidate-authored specifications, including the Bloom progression
    indicator, XP display, and session summary layout.
  \item \textbf{Mapbox integration.} The \texttt{InteractiveMap} component and
    associated Mapbox GL JS configuration.
  \item \textbf{Vercel deployment configuration.} Environment variable setup,
    preview deployment configuration, and build pipeline scaffolding.
  \item \textbf{Prisma schema iteration.} Incremental schema additions from
    candidate-authored specifications, including migration file generation.
  \item \textbf{Seed script scaffolding.} Structure and helper functions for
    the question bank seed; all question content was authored by the candidate.
\end{itemize}

% ─────────────────────────────────────────────────────────────────────────────
\section*{Candidate's Independent Contributions}
\label{app:candidate_contributions}
% ─────────────────────────────────────────────────────────────────────────────

The following constitute the candidate's independent intellectual work and
were not produced by AI tools:

\begin{itemize}[leftmargin=2em]
  \item The hybrid IRT 3PL + BKT algorithm architecture, including the
    selection of EAP over MLE estimation and the full mathematical
    derivation of the credit-factor BKT modification.
  \item The Bloom escalation gating design (0.40 and 0.60 thresholds) and
    its implementation as a computational selection constraint.
  \item The 13-KC curriculum ontology and prerequisite graph, derived from
    the AQA GCSE Geography specification.
  \item The research design: within-subjects quasi-experimental structure,
    VanLehn benchmark selection, 7-day retention protocol, and the Sprint~5
    scope reduction justification.
  \item The five-track multi-method evaluation framework developed in
    response to the $n = 4$ constraint.
  \item All Architecture Decision Records (ADRs) documenting algorithm
    selection rationale.
  \item All Notion sprint planning documentation, retrospective notes, and
    risk register entries.
  \item The IRT calibration strategy: fixing $c = 0.25$, monitoring RMSEA,
    and the two-stage R-to-Node.js parameter pipeline.
\end{itemize}

AI assistance was employed exclusively as an implementation accelerator for
tasks of a non-intellectual, repetitive, or infrastructural nature. All
design decisions, algorithmic contributions, and evaluation methodology
are the candidate's own work and are grounded in the peer-reviewed
literature cited throughout this report.
